\section{How to Read This Report}\label{sec:howtoread}

This report exists in two forms:

\begin{itemize}
  \item \textbf{Slim master} (\texttt{REPLICATION\_EVALUATION\_REPORT\_slim.pdf}):
    aggregate statistics, cross-cutting findings, wave summaries, and a
    top-line per-paper table with hyperlinks to each paper's detailed report.
    This is the default view ($\sim$20 pages).
  \item \textbf{Full master} (\texttt{REPLICATION\_EVALUATION\_REPORT\_full.pdf}):
    everything in the slim master \emph{plus} the complete per-paper evaluation
    subsections inline ($\sim$90 pages).
\end{itemize}

\noindent\textbf{Per-paper detail convention.} Each replication project
maintains a \texttt{REPORT.md} file as its canonical detailed report. The
top-line table in \S\ref{sec:topline} links to these files via relative paths.
To read a paper's full analysis, follow the link (or navigate to
\texttt{<project-dir>/REPORT.md} in the repository).

When a new paper is added to the portfolio:
\begin{enumerate}
  \item Create \texttt{papers/<paper-id>.tex} with the per-paper \LaTeX{} subsection.
  \item Add an entry to \texttt{common/top\_line\_table.tex}.
  \item Add a one-line \verb|\input| to the appropriate wave in
    \texttt{common/wave\_summaries.tex} (inside the
    \verb|\iffull| guard).
  \item Run \texttt{make all} to rebuild both PDFs.
\end{enumerate}
